\documentclass[11pt,a4paper]{report}

\usepackage[margin=1in]{geometry}
\usepackage{listings}
\usepackage{xcolor}
\usepackage{booktabs}
\usepackage{longtable}
\usepackage{hyperref}
\usepackage{caption}
\usepackage{amsmath}
\usepackage{enumitem}
\usepackage{fancyhdr}

\hypersetup{colorlinks=true, linkcolor=blue, urlcolor=blue, citecolor=blue}

\definecolor{codebg}{rgb}{0.96,0.96,0.96}
\definecolor{codekw}{rgb}{0.0,0.0,0.6}
\definecolor{codecomment}{rgb}{0.3,0.5,0.3}
\definecolor{codestring}{rgb}{0.6,0.0,0.0}

\lstdefinestyle{cppstyle}{
  language=C++,
  backgroundcolor=\color{codebg},
  basicstyle=\ttfamily\small,
  commentstyle=\color{codecomment}\itshape,
  keywordstyle=\color{codekw}\bfseries,
  stringstyle=\color{codestring},
  numbers=left,
  numberstyle=\tiny\color{gray},
  numbersep=8pt,
  frame=single,
  rulecolor=\color{gray!40},
  breaklines=true,
  breakatwhitespace=false,
  showstringspaces=false,
  tabsize=4,
  columns=fullflexible,
  keepspaces=true
}
\lstset{style=cppstyle}

\pagestyle{fancy}
\fancyhf{}
\fancyhead[L]{APARA Simulator (engine\_isp) --- Source Fixes: What Changed and Why}
\fancyhead[R]{\thepage}
\renewcommand{\headrulewidth}{0.4pt}

\title{\textbf{APARA Simulator (\texttt{engine\_isp}) \\ Source Fixes Applied During Compiler Verification \\ \large What Changed, and Why Each Change Was Necessary}}
\author{APARA C Compiler Project --- Verification \& Audit}
\date{\today}

\begin{document}

\maketitle
\tableofcontents

% ============================================================================
\chapter{Introduction: Why These Edits Were Made}
% ============================================================================

\section{Context}

The APARA C compiler is now feature-complete (integer C, floating point,
function pointers, variadic functions, $>$4 arguments). To find any remaining
defects before hardware bring-up, we ran a \textbf{differential verification
campaign}: 1{,}013 C programs (a 13-test directed battery plus 1{,}000
randomly generated whole programs), each compiled twice --- once by the APARA
compiler and executed on the \texttt{engine\_isp} simulator, and once natively
by \texttt{gcc}. Every \texttt{results[]} slot must match \texttt{gcc}
bit-for-bit. The generated programs are constructed to be free of undefined
behaviour, so \emph{every mismatch is a genuine toolchain defect}, either in
the compiler or in the simulator.

A coverage auditor additionally proves that the campaign leaves nothing
untouched: all 102 (instruction $\times$ sub-operation $\times$ type-tag)
combinations the compiler can emit are exercised by at least one passing,
\texttt{gcc}-verified program.

\section{Why the simulator source had to be edited}

Six defects found by this campaign are in the \emph{simulator}, not the
compiler. They could not be worked around from the compiler side:

\begin{itemize}[itemsep=2pt]
  \item Three make \textbf{ISA-documented instructions produce silently wrong
        values} (\texttt{\$fsqrt} returns 0; unsigned \texttt{\$vreduce}
        treats lanes as signed; scalar float$\to$int casts lose their sign
        extension). A compiler cannot emit correct code for an instruction the
        simulator executes incorrectly --- and ``silently wrong'' is the worst
        failure mode for hardware validation, because the RTL would be judged
        against a wrong reference.
  \item Three are \textbf{re-applications of the July 17--18 floating-point
        fixes} (cast dispatch swap, the unimplemented float-cast branch, and
        the \texttt{cast\_int\_to\_float} promotion bug). Those had been fixed
        and deployed as binaries during the FP campaign, but the source tree
        had been left pristine; a later rebuild silently reverted them
        (Chapter~\ref{ch:process}). They are now fixed \emph{in source}.
\end{itemize}

Each edit is \textbf{minimal and local}: it uses machinery that already exists
in the file (e.g.\ \texttt{\$fsqrt} is routed through the existing ALU
dispatch, which already implements \texttt{fp\_sqrt}), changes no interfaces,
and alters no behaviour outside the defective case. After \emph{every} rebuild
the full regression gate was re-run (feature\_sweep 16/16, universal 21/21,
pointer\_bugs 15/15, fp01--09 50/50), and the final campaign result is
\textbf{613 PASS / 0 FAIL / 0 HANG} with 102/102 coverage.

\section{Summary of all changes}

\begin{longtable}{@{}p{0.5cm}p{3.6cm}p{4.3cm}p{5.6cm}@{}}
\toprule
\# & File / function & Defect & Why the edit was necessary \\
\midrule
1 & \texttt{MachineRun.cpp} \newline \texttt{McodeFsqrtInstruction::Execute} &
    \texttt{\$fsqrt} was an unimplemented stub returning 0 &
    \texttt{\$fsqrt} is in the ISA and the assembler grammar; every
    \texttt{sqrt}/\texttt{sqrtf} silently produced 0 \\
\addlinespace
2 & \texttt{McodeExecute.cpp} \newline \texttt{\_\_\_execute\_cast\_operation\_\_\_} &
    float casts fell into an error branch, output register left
    \emph{uninitialized} &
    all f32/f64 $\leftrightarrow$ int casts landed host-memory garbage in
    registers (re-applied Jul-17 fix) \\
\addlinespace
3 & \texttt{McodeOperations.cpp} \newline \texttt{\_\_\_cast\_operation\_\_\_} &
    int$\leftrightarrow$float dispatch called the two helpers swapped &
    every conversion read its operand as the wrong representation
    (re-applied Jul-17 fix) \\
\addlinespace
4 & \texttt{McodeOperations.cpp} \newline \texttt{\_\_\_cast\_operation\_\_\_} &
    scalar 32-bit casts misclassified as \emph{vector} casts &
    the result was masked to 32 bits, destroying the sign extension of
    negative float$\to$int results \\
\addlinespace
5 & \texttt{McodeOperations.cpp} \newline \texttt{\_\_vreduce\_operation\_\_} &
    unsigned \texttt{\$vreduce} sign-extended its lanes &
    \texttt{\$vu8/\$vu16/\$vu32} reductions summed/compared lanes as negative
    numbers \\
\addlinespace
6 & \texttt{McodeFpuUtils.cpp} \newline \texttt{cast\_int\_to\_float} &
    C ternary promotion turned negative ints into $2^{64}-n$ &
    every \emph{negative} int$\to$float cast was wrong
    (re-applied Jul-18 fix) \\
\bottomrule
\end{longtable}

% ============================================================================
\chapter{The Six Changes in Detail}
% ============================================================================

\section{\texttt{\$fsqrt} executed as an unimplemented stub (NEW)}

\textbf{File / function:} \texttt{MachineRun.cpp},
\texttt{McodeFsqrtInstruction::Execute}

\textbf{Symptom.} Every \texttt{\$fsqrt} produced 0:
\texttt{(long long)sqrt(4.0)} returned 0 instead of 2. Found by directed test
\texttt{d08\_fsqrt\_float.c}; the simulator printed
``\texttt{McodeFsqrtInstruction yet to be implemented}'' and set the output
to 0.

\textbf{Root cause.} The \texttt{Execute} body was a placeholder. All the
machinery already existed: the constructor sets opcode \texttt{\_\_FSQRT},
and \texttt{\_\_alu\_operation\_\_} already dispatches \texttt{\_\_FSQRT} to
\texttt{fp\_sqrt} --- the instruction class simply never invoked it.

\textbf{Why we edited it.} \texttt{\$fsqrt} is a documented ISA instruction
that the compiler lowers \texttt{sqrt()}/\texttt{sqrtf()} to; without this
fix the instruction cannot be validated at all, and a hardware implementation
would have no correct simulator reference. The edit only \emph{connects} the
stub to the existing, already-tested ALU execution path --- no new arithmetic
code was written.

\textbf{Before:}
\begin{lstlisting}
void McodeFsqrtInstruction::Execute     (McodeMachine* mc)
{
    McodeRoot::Error("McodeFsqrtInstruction yet to be implemented", this);
    this->McodeInstruction::Set_Out_Arg (0, 0);
    return;
}
\end{lstlisting}

\textbf{After:}
\begin{lstlisting}
void McodeFsqrtInstruction::Execute     (McodeMachine* mc)
{
    uint64_t ovalue;
    ___execute_alu_operation___ (this->McodeInstruction::Get_Opcode(),
            this->Get_Dest_Type(),
            this->Get_Src_Type(),
            this->McodeInstruction::Get_In_Arg(0),
            0, 0, 0,
            ovalue);
    this->McodeInstruction::Set_Out_Arg (0, ovalue);
    if(__global_verbose_flag)
        McodeRoot::Info(Int64ToString (this->McodeInstruction::Get_Address()) + ": Fsqrt!");
    return;
}
\end{lstlisting}

\textbf{Verification.} \texttt{sqrt(4.0)}, \texttt{sqrt(2.25)},
\texttt{sqrtf(9.0f)} etc.\ all bit-match gcc (d08 = 8/8 golden checks);
run-log shows \texttt{0x4010000000000000 $\to$ 0x4000000000000000}
($4.0 \to 2.0$, exact).

% ----------------------------------------------------------------------------
\section{Float casts fell into an unimplemented error branch (RE-APPLIED)}

\textbf{File / function:} \texttt{McodeExecute.cpp},
\texttt{\_\_\_execute\_cast\_operation\_\_\_}

\textbf{Symptom.} Any \texttt{\$cast} with a float type on either side printed
``\texttt{Float conversions in cast not supported as yet.}'' and left the
output array \emph{uninitialized} --- host-process stack garbage (values such
as \texttt{0x7ffd83e474a0}) landed in the destination register.

\textbf{Root cause.} The complete implementation (which forwards to
\texttt{\_\_\_cast\_operation\_\_\_}, and that function dispatches
int$\leftrightarrow$float itself) was guarded by
\texttt{if (neither type is float)}; float casts took the error branch.

\textbf{Why we edited it.} Floating-point support in the compiler is verified
bit-exactly against gcc (fp01--fp09), and every float$\leftrightarrow$int
conversion goes through \texttt{\$cast}. Reading uninitialized memory is also
nondeterministic --- results changed from run to run, which would make any
hardware-vs-simulator comparison meaningless. The guard was removed so every
cast takes the one real implementation; no conversion logic was added here.

\textbf{Before:}
\begin{lstlisting}
    if(!dest_type.Get_Float_Flag()
            && !src_type.Get_Float_Flag())
    {
        /* ... full implementation ... */
    }
    else
    {
        McodeRoot::Error ("Float conversions in cast not supported as yet.", NULL);
    }
\end{lstlisting}

\textbf{After:}
\begin{lstlisting}
    if(1)   /* every cast takes the main path; ___cast_operation___
               dispatches int<->float itself */
    {
        /* ... full implementation (unchanged) ... */
    }
\end{lstlisting}

\textbf{Verification.} fp01--fp09 (50/50 golden checks, bit-exact incl.\
rounding); campaign-wide float casts in 600+ random programs.

% ----------------------------------------------------------------------------
\section{int$\leftrightarrow$float cast dispatch called the helpers swapped (RE-APPLIED)}

\textbf{File / function:} \texttt{McodeOperations.cpp},
\texttt{\_\_\_cast\_operation\_\_\_}

\textbf{Symptom.} Every int$\to$float and float$\to$int conversion produced
garbage: each helper interpreted its operand as the wrong representation
(integer bits read as IEEE and vice versa).

\textbf{Root cause.} The dispatch was inverted: an \emph{integer source} was
sent to \texttt{cast\_float\_to\_int}, and a float source with integer
destination to \texttt{cast\_int\_to\_float}.

\textbf{Why we edited it.} This is a plain transposition bug --- the two
correctly implemented helpers were wired to the opposite conditions. The edit
swaps the two call sites back; nothing else changes.

\textbf{Before:}
\begin{lstlisting}
            if(!src_type.Get_Float_Flag())
                w = cast_float_to_int (src_type, dest_type, v);
            else if(!dest_type.Get_Float_Flag())
                w = cast_int_to_float (src_type, dest_type, v);
            else
                w = cast_float_to_float (src_type, dest_type, v);
\end{lstlisting}

\textbf{After:}
\begin{lstlisting}
            if(!src_type.Get_Float_Flag())
                w = cast_int_to_float (src_type, dest_type, v);
            else if(!dest_type.Get_Float_Flag())
                w = cast_float_to_int (src_type, dest_type, v);
            else
                w = cast_float_to_float (src_type, dest_type, v);
\end{lstlisting}

\textbf{Verification.} Same as change 2 (all casts flow through here).

% ----------------------------------------------------------------------------
\section{Scalar 32-bit casts misclassified as vector casts (NEW)}

\textbf{File / function:} \texttt{McodeOperations.cpp},
\texttt{\_\_\_cast\_operation\_\_\_}

\textbf{Symptom.} \texttt{(int)(-6.0f)} stored \texttt{0x00000000FFFFFFFA}
instead of the sign-extended \texttt{0xFFFFFFFFFFFFFFFA}; regression test
fp03 caught it. The conversion itself was correct
(\texttt{cast\_float\_to\_int} returns the full sign-extended 64-bit value)
--- the result was then masked back to 32 bits.

\textbf{Root cause.} Vector-ness was inferred from the number of unpacked
input elements, \texttt{in\_vals.size() > 1}. But \texttt{Break\_Vector}
splits even a \emph{scalar} 32-bit source register into two 32-bit entries,
so every scalar cast with a 32-bit source was treated as a vector cast, and
vector casts mask each output to the destination lane width.

\textbf{Why we edited it.} The type words themselves carry an explicit
\texttt{Vector\_Flag}; inferring vector-ness from a size heuristic contradicts
the instruction's own declared type and cannot distinguish ``scalar f32''
from ``vector of two f32 lanes''. Using the declared flags makes the scalar
path take the unmasked (sign-preserving) result, and leaves true vector casts
exactly as before.

\textbf{Before:}
\begin{lstlisting}
    int is_vector_cast = (in_vals.size() > 1);
\end{lstlisting}

\textbf{After:}
\begin{lstlisting}
    int is_vector_cast = (src_type.Get_Vector_Flag() && dest_type.Get_Vector_Flag());
\end{lstlisting}

\textbf{Verification.} fp03 (negative float$\to$int) restored to 4/4;
full gate 61/61; d02 cast matrix 8/8.

% ----------------------------------------------------------------------------
\section{Unsigned \texttt{\$vreduce} sign-extended its lanes (NEW)}

\textbf{File / function:} \texttt{McodeOperations.cpp},
\texttt{\_\_vreduce\_operation\_\_}

\textbf{Symptom.} \texttt{\_\_vreduce\_vu8(0x01f2030405060708)} returned
\texttt{0x14} (20) instead of \texttt{0x114} (276): the \texttt{0xf2} lane
was summed as $-14$ instead of $242$. \texttt{\_\_vreduce\_max\_vu8} likewise
picked the wrong maximum. Found by directed test \texttt{d06\_vreduce.c}.

\textbf{Root cause.} A shared lane variable \texttt{r} was computed with
\texttt{Sign\_Extend\_64} \emph{before} the signed/unsigned branch; the
signed branch shadows it with its own (correct) sign-extended copy, but the
\emph{unsigned} branch used the outer, sign-extended value.

\textbf{Why we edited it.} The entire point of the \texttt{\$vu*} type tags is
zero-extension of lanes; with this bug the unsigned reductions were
indistinguishable from the signed ones, i.e.\ half of the
\texttt{\$vreduce} ISA surface was unimplemented in effect. The lane value
arriving from \texttt{Break\_Vector} is already masked to the lane width, so
the unsigned branch simply uses it directly. The signed branch is untouched.

\textbf{Before:}
\begin{lstlisting}
            // sign extend to 64-bits.
            int64_t r = (int64_t) Sign_Extend_64(src_type.Get_Nbits()-1, ele);
            if(signed_flag)
            {
                int64_t r = (int64_t) Sign_Extend_64(...);  // shadows outer r
                ...
            }
            else
            {
                ... uses the OUTER, sign-extended r ...
            }
\end{lstlisting}

\textbf{After:}
\begin{lstlisting}
            uint64_t r = ele;      // zero-extended: ele is already lane-masked
            if(signed_flag)
            {
                int64_t r = (int64_t) Sign_Extend_64(...);  // unchanged
                ...
            }
            else
            {
                ... unsigned branch now sums/compares zero-extended lanes ...
            }
\end{lstlisting}

\textbf{Verification.} d06 = 8/8 golden checks covering
\texttt{+}/\texttt{\$max} at \texttt{vi8/vu8/vi16/vu16/vi32/vu32} vs gcc.

% ----------------------------------------------------------------------------
\section{\texttt{cast\_int\_to\_float}: ternary promotion of negatives (RE-APPLIED)}

\textbf{File / function:} \texttt{McodeFpuUtils.cpp},
\texttt{cast\_int\_to\_float}

\textbf{Symptom.} \texttt{(float)(-3)} became $2^{64}-3 \approx 1.8\times
10^{19}$ as a double, then saturated on any conversion back. All
\emph{positive} int$\to$float casts were unaffected, which is how the bug
originally escaped notice.

\textbf{Root cause.} In
\texttt{cond ? (u \& 0xffffffff) : (int) u}
the C conditional operator promotes both arms to a common type --- here the
unsigned arm's \texttt{uint64\_t} --- so the signed arm's negative
\texttt{int} value was converted back to a huge unsigned number before
becoming a double.

\textbf{Why we edited it.} This is a classic C promotion pitfall, invisible
in code review but fatal for every negative operand. The replacement makes
the signedness handling explicit (branch, then sign-extend from the actual
source width) --- it also fixes the old unsigned arm's hard-coded 32-bit
mask, which mishandled unsigned 8/16/64-bit sources.

\textbf{Before:}
\begin{lstlisting}
    double x = (src_type.Get_Unsigned_Flag() ?  (u & 0xffffffff) : (int) u);
\end{lstlisting}

\textbf{After:}
\begin{lstlisting}
    double x;
    if(src_type.Get_Unsigned_Flag())
    {
        x = (double) (u & __mmask__(snbits));
    }
    else
    {
        int64_t sv = ((int64_t) (u << (64 - snbits))) >> (64 - snbits);
        x = (double) sv;
    }
\end{lstlisting}

\textbf{Verification.} fp03 (positive and negative int$\to$float, $\times$,
$\div$, $+$) = 4/4 vs gcc; hand-written \texttt{\$cast} micro-test from the
original Jul-18 session.

% ============================================================================
\chapter{Checked and Found Already Correct (No Change Made)}
% ============================================================================

For completeness, the campaign also re-audited areas with a history of
defects; the following are \textbf{correct in the current source} and were
left untouched:

\begin{itemize}
  \item \texttt{McodeFpu.hpp} \texttt{\_\_float32\_sub\_\_} /
        \texttt{\_\_float64\_sub\_\_}: the July~17 \texttt{fp\_sub} defects
        (bit-pattern subtraction at 32-bit; `$+$' instead of `$-$' at 64-bit)
        are \emph{not present} --- both macros convert through double and
        subtract correctly.
  \item The unsigned scalar ALU path (\texttt{\_\_uexec\_64\_\_} on raw
        \texttt{uint64\_t} + \texttt{CastToU64}) is correct, including
        unsigned divide and logical right shift. The older build's
        \texttt{\$u64}-zeroing defect no longer exists; consequently the
        compiler now emits \texttt{(\$u64)} for unsigned \texttt{/},
        \texttt{\%} (its divide step) and \texttt{>>}, all verified against
        gcc (including a logical shift by 63).
  \item The July~5 backward-\texttt{\$call} sign-extension fix is present;
        every call to an earlier-defined function in the 61-test gate is a
        backward call.
\end{itemize}

Two \emph{behavioural} findings in the assembler worth knowing (not changed):

\begin{itemize}
  \item \texttt{\$pack} requires \texttt{packed\_nbits \% word\_nbits == 0}
        (\texttt{McodeClasses.hpp:412}); violations are reported only as the
        generic ``\texttt{there were invalid instructions}'' with no
        per-instruction message.
  \item \texttt{mcode\_align} reports lane-placement failures ($>$4
        load/store or $>$1 divide/sqrt per bundle) on stderr but \emph{still
        emits the bundle} with the control-transfer instruction not moved to
        lane~0 --- downstream this caused a return into the middle of a
        bundle. The compiler's bundler now enforces these limits itself, and
        our scripts no longer discard align/assemble stderr.
\end{itemize}

% ============================================================================
\chapter{Process Note: Keeping Fixes in Source}
\label{ch:process}
% ============================================================================

Changes 2, 3 and 6 were originally made during the July 17--18
floating-point campaign, but only the \emph{binaries} were deployed; the
source tree was left pristine. When the toolchain was rebuilt on July~19
(for the \texttt{\$fsqrt} fix), the rebuild regenerated all binaries from the
unfixed source and silently reverted those three fixes. The regression gate
caught this immediately (fp03 failing, 60/61), and all fixes were re-applied
\emph{in source} with explanatory comments.

Two safeguards are now in place:

\begin{itemize}
  \item \texttt{testing/fuzz1000/check\_engine\_fixes.sh} greps the simulator
        source for all six fixes; it is run before and after any rebuild.
  \item The full regression gate (61 golden tests) plus, when relevant, the
        fuzz campaigns are re-run after \emph{every} toolchain rebuild.
\end{itemize}

\textbf{Recommendation for the upstream tree:} absorbing the six diffs in
Chapter~2 brings the upstream simulator in line with the reference against
which the compiler is verified. Each is minimal, local, and documented at the
change site with a comment beginning ``\texttt{Re-applied fix}'' or referring
to the fuzz1000 campaign.

\end{document}
